% !TEX root = begin.tex
% !TeX spellcheck = nl_NL
% \titelpagina definieert de titlepage
% 1 = Naam van het plaatje in directory ./fig. 
%     Neem een hoog resolutie plaatje.
%     De hoogte moet ongeveer 0.5 x de breedte zijn.
%>titelpagina
%\titelpagina{logo-hr.png}
%<titelpagina
% commando voor versiehistory tabel
% #1 = tabel inhoud
%      tabel heeft 4 kolommen:
%      Datum & Versie & Omschrijving & Auteur
%      Gebruik \\ \line om de rijen te scheiden.
%>versiehistorie
\versiehistorie{
	24-03+2026 & \huidigeversie &
	Driefasen toegevoegt & JHT
	\\ \hline
	24-03-2026 & 0.2 &
	Hoofdstukken 1, 2 en 3 toegevoegt
	& JHT
	\\ \hline
    24-03-2026 & 0.1 & 
	Instellingen, versiehistory en titelpagina toegevoegt 
	& JHT 
}

%<versiehistorie
% \footnotetext{Toelichting versiecodering $A$.$Bc$: $A$ = grote aanpassing, $B$ = kleine aanpassing, $c$ = taal- of wiskundige
% correcties.}

% In principe is het onzinnig om versiehistorie in het document zelf bij te houden als de broncode beschikbaar is in een publiek toegankelijke VCS (version control system). 
% Maar omdat dit niet voor alle documenten het geval hoeft te zijn, wordt hier toch, als voorbeeld, een versiehistorietabel gebruikt.\noclub[4]

%>licentie    
\vspace*{\stretch{100}}
%\includegraphics{figs/by-nc-sa_eu.pdf}
\par
%\documenttype{} \titel{} van Hogeschool Rotterdam is in licentie gegeven volgens een \href{http://creativecommons.org/licenses/by-nc-sa/3.0/nl/} {Creative Commons Naamsvermelding\hyp{}NietCommercieel\hyp{}GelijkDelen 3.0 Nederland\hyp{}licentie}.
%<licentie
%\Cref{fig:zebra,kleinezebra,fig:binary,fig:fractal} in deze \MakeLowercase\documenttype{} zijn afkomstig van \href{http://pixabay.com}{Pixabay} en vallen onder het publieke domein volgens een 
%\href{http://creativecommons.org/publicdomain/zero/1.0/deed.nl}{Publiek Domein Verklaring (CC0 1.0)}.

% De \LaTeX\hyp{}code van deze \MakeLowercase\documenttype{} kun je vinden op: \url{https://bitbucket.org/HR_ELEKTRO/latex-stijl}.